% ============================================================
%  ROB 311 H1  –  KR & Probabilistic Inference  (Harder Variants)
% ============================================================
%
%  Assumes the following in includes.tex (or preamble):
%    \usepackage{enumerate, enumitem, amsmath, amssymb, tikz, xcolor, booktabs}
%    \newcounter{qnum}
%    \newcommand{\question}[2]{...}        % as before
%    \newcommand{\points}[1]{[#1 pts]}
%    \newenvironment{sol}{...}{...}         % blue solution block
%    \newcommand{\choice}[1]{\underline{\hspace{1cm}}}  % circle/write answer
%    \newcommand{\numbox}[1]{\framebox[2cm]{\rule{0pt}{10pt}}}
%
%  \choice produces a labelled blank for multiple-choice items.
%  \numbox produces a box for a numerical answer.
%  Students are expected to show work on the page; final answers
%  in the boxes allow rapid grading.
% ============================================================

\newcounter{qnum}
\setcounter{qnum}{0}
\newcommand{\question}[2]{%
  \stepcounter{qnum}%
  \noindent\textbf{Question \theqnum} \hfill \textbf{[#1 pts]}\\[1mm]
  \noindent #2
  \vspace{4mm}
}


% ============================================================
%  QUESTION 2  –  Clausal Form Conversion  (new)
% ============================================================
\question{12}{%
\textbf{Conversion to Clausal Form.}\\[2mm]
Convert the following FOL sentence to clausal (conjunctive normal) form, showing \textbf{each transformation step} and naming it. Write your final clauses in the box provided.

\[
  \varphi \;=\; \forall x\,\Bigl(
    \texttt{animal}(x)
    \;\to\;
    \bigl(\exists y\,(\texttt{food}(y) \wedge \texttt{eats}(x,y))\bigr)
    \;\vee\;
    \neg\,\forall z\,(\texttt{plant}(z) \to \texttt{eats}(x,z))
  \Bigr)
\]

\noindent The steps you must show, in order, are:
\begin{enumerate}
  \item Eliminate implications.
  \item Move negations inward (NNF).
  \item Standardise variables apart.
  \item Skolemise (eliminate existentials).
  \item Drop universal quantifiers.
  \item Distribute $\vee$ over $\wedge$ (CNF).
\end{enumerate}

\vspace{4mm}
\noindent\textbf{Show each step below.} After every step, write the resulting formula. If a step is not required, you may write \texttt{N/A} for it.

\vspace{3mm}
\noindent\textbf{Step 1 – Eliminate implications.}
\vspace{2.5cm}

\noindent\textbf{Step 2 – Move negations inward.}
\vspace{2.5cm}

\noindent\textbf{Step 3 – Standardise variables apart.}
\vspace{2.5cm}

\clearpage
\noindent\textbf{Step 4 – Skolemise} (indicate any Skolem constants/functions).
\vspace{2.5cm}

\noindent\textbf{Step 5 – Drop universal quantifiers.}
\vspace{2.5cm}

\noindent\textbf{Step 6 – Distribute $\wedge$ over $\vee$.}
\vspace{2.5cm}

\noindent\textbf{Final clauses} (write each clause in its own box - you may not need them all): \\~\\
\noindent\numbox[14cm]\\
\numbox[14cm]\\
\numbox[14cm]\\
\numbox[14cm]\\
\numbox[14cm]\\
\numbox[14cm]\\
}

\clearpage

% ============================================================
%  QUESTION 3  –  Logical Ambiguity + Structures  (harder)
% ============================================================
\question{12}{%
\textbf{Logical Ambiguity and Structures.}\\[2mm]
Consider the following knowledge base $\Phi$ in clausal form:
\begin{enumerate}[(1)]
  \item $\texttt{raptor}(\texttt{a})$
  \item $\texttt{nocturnal}(\texttt{b})$
  \item $\neg\,\texttt{nocturnal}(x) \;\vee\; \texttt{raptor}(x)$
  \item $\neg\,\texttt{raptor}(x) \;\vee\; \texttt{hunts}(x, \texttt{mouse})$
  \item $\neg\,\texttt{raptor}(x) \;\vee\; \texttt{nocturnal}(x) \;\vee\; \texttt{diurnal}(x)$
  \item $\neg\,\texttt{nocturnal}(x) \;\vee\; \neg\,\texttt{diurnal}(x)$
\end{enumerate}
Let $\Phi' = \texttt{nocturnal}(\texttt{a})$.

\begin{enumerate}[(a)]

\item \points{2} Is $\Phi'$ logically ambiguous given $\Phi$? Circle: \;\choice{Yes} \;/\; \choice{No}\\[2mm]
Justify in 1--2 sentences.

\item \points{5} Construct a structure $\mathcal{S}_1$ such that $\mathcal{S}_1 \vDash \Phi$ \textbf{and} $\mathcal{S}_1 \vDash \Phi'$.\\
Specify domain, constant interpretations, and predicate interpretations.
Verify \textbf{every} clause.


\item \points{5} Construct a structure $\mathcal{S}_2$ such that $\mathcal{S}_2 \vDash \Phi$ \textbf{and} $\mathcal{S}_2 \not\vDash \Phi'$.\\
Specify domain, constant interpretations, and predicate interpretations.
Verify \textbf{every} clause.


\end{enumerate}
}

\clearpage

% ============================================================
%  QUESTION 4  –  Resolution by Refutation  (harder)
% ============================================================
\question{12}{%
\textbf{Resolution by Refutation.}\\[2mm]
Consider the knowledge base $\Phi$ in clausal form:
\begin{enumerate}
  \item $\texttt{scout}(\texttt{c0})$
  \item $\texttt{forager}(\texttt{c0})$
  \item $\texttt{target}(\texttt{r1})$
  \item $\neg\,\texttt{scout}(x) \;\vee\; \neg\,\texttt{target}(y) \;\vee\; \texttt{marks}(x, y)$
  \item $\neg\,\texttt{forager}(x) \;\vee\; \neg\,\texttt{marks}(x, y) \;\vee\; \texttt{exploits}(x, y)$
  \item $\neg\,\texttt{exploits}(x, y) \;\vee\; \texttt{depletes}(y)$
  \item $\neg\,\texttt{depletes}(y) \;\vee\; \neg\,\texttt{target}(y) \;\vee\; \texttt{scarce}(y)$
\end{enumerate}
We wish to prove $\Phi' = \texttt{scarce}(\texttt{r1})$ by \textbf{resolution by refutation}.

\begin{enumerate}

\item \points[2] Write the clause that must be added to begin the proof.

\begin{enumerate}
    \item[8.] \numbox[12cm]
\end{enumerate}


\item \points[8] Carry out the full resolution proof labeling each step using the notation \texttt{\#.\;r[i,j]\{$\delta$\}}, where \texttt{\#.}\;is the new clause number, \texttt{i} and \texttt{j} are the clauses being resolved, and $\delta$ is a substitution if required.


\end{enumerate}
}

\clearpage

\begin{center}\textit{(extra space)}\end{center}
\clearpage

% ============================================================
%  QUESTION 5  –  D-Separation  (harder)
% ============================================================
\question{10}{%
\textbf{Bayesian Network: D-Separation.}\\[2mm]
Consider the Bayesian network below, which models a wildfire detection system:

\vspace{3mm}
\begin{center}
\begin{tikzpicture}[
    > = stealth, shorten > = 1pt, auto, node distance = 2.2cm, semithick
]
\tikzstyle{every state}=[draw=black, thick, fill=white, minimum size=7mm, font=\small]

\node[state] (L) {$L$};
\node[state] (H) [right of=L] {$H$};
\node[state] (W) [right of=H] {$W$};
\node[state] (F) [below of=H] {$F$};
\node[state] (S) [below left of=F] {$S$};
\node[state] (A) [below right of=F] {$A$};
\node[state] (R) [below of=A] {$R$};

\path[->] (L) edge (F);
\path[->] (H) edge (F);
\path[->] (W) edge (F);
\path[->] (W) edge (A);
\path[->] (F) edge (S);
\path[->] (F) edge (A);
\path[->] (A) edge (R);

\node[above=2mm of L, font=\footnotesize\color{gray}] {Lightning};
\node[above=2mm of H, font=\footnotesize\color{gray}] {Heat};
\node[above=2mm of W, font=\footnotesize\color{gray}] {Wind};
\node[left=2mm of S, font=\footnotesize\color{gray}] {Smoke};
\node[right=2mm of A, font=\footnotesize\color{gray}] {Alarm};
\node[right=2mm of R, font=\footnotesize\color{gray}] {Response};
\node[left=1mm of F, font=\footnotesize\color{gray}] {Fire};
\end{tikzpicture}
\end{center}

\noindent For each query, circle \textbf{one} answer and provide a justification.

\vspace{3mm}
\begin{enumerate}

  \item \points[2] $L$ is unconditionally independent of $W$ \\
  \hfill \choice[True] \choice[False] \choice[Cannot be determined]
\clearpage
  \item \points[2] $L$ is independent of $W$ given $F$ \\
  \hfill \choice[True] \choice[False] \choice[Cannot be determined]
\vspace{5cm}
  \item \points[3] $S$ is independent of $R$ given $F$ \\
  \hfill \choice[True] \choice[False] \choice[Cannot be determined]
\vspace{5cm}
  \item \points[3] $H$ is independent of $W$ given $A$ \\
  \hfill \choice[True] \choice[False] \choice[Cannot be determined]

\end{enumerate}
}

\clearpage

% ============================================================
%  QUESTION 6  –  Probabilistic Inference with Scaffolding
% ============================================================
\question{16}{%
\textbf{Bayesian Network: Variable Elimination with Scaffolding.}\\[2mm]
Consider a Bayesian network with binary variables $E$ (Earthquake), $B$ (Burglar), $A$ (Alarm), where $E \to A$, $B \to A$ (i.e.\ $A$ has parents $E$ and $B$, and $E$, $B$ are marginally independent).

\vspace{3mm}
\begin{center}
\begin{tabular}{|c|c|}
\hline $E$ & $P(E)$ \\ \hline
$0$ & $0.98$ \\ $1$ & $0.02$ \\ \hline
\end{tabular}
\hspace{6mm}
\begin{tabular}{|c|c|}
\hline $B$ & $P(B)$ \\ \hline
$0$ & $0.99$ \\ $1$ & $0.01$ \\ \hline
\end{tabular}
\hspace{6mm}
\begin{tabular}{|c|c|c|c|}
\hline $E$ & $B$ & $A$ & $P(A \mid E, B)$ \\ \hline
$0$ & $0$ & $1$ & $0.001$ \\
$0$ & $1$ & $1$ & $0.94$ \\
$1$ & $0$ & $1$ & $0.29$ \\
$1$ & $1$ & $1$ & $0.95$ \\ \hline
\end{tabular}
\end{center}

\vspace{4mm}
\noindent We wish to compute $P(B{=}1 \mid A{=}1)$.
\begin{enumerate}
    \item \points[2] Write $\textsf{pr}(B=1|A=1)$ in terms of the joint distribution between $A$ and $B$ (i.e., $\textsf{pr}(A,B)$).
\[\textsf{pr}(B=1|A=1) = \numbox[5cm] \; \bigg/ \; \numbox[5cm]\]

\item \points[2] Write $\textsf{pr}(A,B)$ in terms of the joint distribution between $A,B$ and $E$ (i.e., $\textsf{pr}(A,B,E)$).
\[\textsf{pr}(A,B) = \numbox[8cm]\]

\item \points[2] Write $\textsf{pr}(A,B,E)$ in terms of the original factors from the Bayesian network. Recall that for a variable, $X$, the associated factor will have the form $\textsf{pr}(X|\textsf{pnt}(X))$, where $\textsf{pnt}(X)$ is the set of $X$'s parents.
\[\textsf{pr}(A,B,E) = \numbox[12cm]\]

\item  \points[4] Calculate $\textsf{pr}(B=1|A=1)$. Show your work.
\[\textsf{pr}(B=1|A=1) = \numbox[5cm]\]
\end{enumerate}
}

\clearpage
\begin{center}\textit{(extra space)}\end{center}
\clearpage